\author{Имя Фамилия}
\doi{}
\authorinfo{Имя Отчество Фамилия}{Аффилиация}{email}

\chapter[Полное название статьи][Короткое название статьи...]{Полное название статьи}

\begin{abstract}

Аннотация

\end{abstract}

\keywords{ключевые слова}



  \subsection{Введение}

Текст\footnote{Текст сноски.} 

\subsection{Основная часть}

Текст [Фамилия год, с. 1--2] 


% Пример оформления таблицы, таблица оформляется в книжном стиле.

\begin{table}[]
\begin{tabular}{@{}ll@{}}
\toprule
Заголовок колонки 1   & Заголовок колонки 2 \\ \midrule
Данные колонки 1      & Данные колонки 2    \\ \bottomrule
\end{tabular}
\caption{Подпись к таблице}
\label{tab:table1}
\end{table}

Ссылка на таблицу в тексте: таблица \ref{tab:table1}

% Пример оформления рисунка

\begin{figure}
	\centering
	\includegraphics[width=1\linewidth]{images/filename}
	\caption{Подпись к рисунку}
	\label{fig:1}
\end{figure}

Ссылка на рисунок в тексте: Рис. \ref{fig:1}

% Пример оформления библиографии

\begin{thebibliography}{99}

\sources

% Список источников


\refs

% Список цитированной литературы

\bibitem[Фамилия год]{} Фамилия\,И. Название. Город: Издательство, год.

\bibitem[Фамилия год]{} Фамилия\,И.\,О. Название // Название журнала. год. №\,1. С.\,1--2.

\end{thebibliography}

